\documentclass[12pt]{article}
\usepackage[utf8]{inputenc}
\usepackage[T1]{fontenc}
\usepackage[margin=1in]{geometry}
\usepackage{setspace}
\usepackage{xcolor}
\usepackage[colorlinks=true, urlcolor=blue!60!black, linkcolor=blue!60!black]{hyperref}
\setlength{\parindent}{0em}
\setlength{\parskip}{0.45em}
\renewcommand{\baselinestretch}{1.0}
\begin{document}
\begin{center}
{\Large \textit{Moore v. City of East Cleveland}, 431 U.S. 494 (1977)\par}
\vspace{0.3em}
{\normalsize Case Brief --- One-Page Edition\par}
\end{center}

\textbf{Memory Jogger.} A city may not make a crime of a grandmother living with her grandsons: extended-family living is deeply rooted in the Nation's history and tradition (\href{https://supreme.justia.com/cases/federal/us/431/494/}{Justia}).

\textbf{Facts.} East Cleveland's housing code defined ``family'' to allow at most one dependent child's children per household. Inez Moore lived with her son and two grandsons who were first cousins; one had lived with her since his mother died. The city ordered her to remove him as an ``illegal occupant'' and prosecuted her when she refused.

\textbf{Procedural History.} Convicted (five days in jail and a \$25 fine); affirmed under \textit{Belle Terre}; Ohio Supreme Court dismissed.

\textbf{Votes.} Reversed, 5--4: Powell for a plurality of four; Brennan concurred; Stevens concurred in the judgment; Burger, Stewart (with Rehnquist), and White dissented.

\textbf{Holding.} Criminalizing cohabitation of close relatives violates the Fourteenth Amendment. The plurality applied close scrutiny to intrusions on the family; Stevens found an arbitrary taking of the owner's right to use her home.

\textbf{Reasoning.} Unlike \textit{Belle Terre}'s limit on unrelated persons, this law intruded on the family itself, with only a ``tenuous'' relation to crowding, traffic, or school costs: it allowed a dozen children of one son but not two cousins, and floor-area limits already controlled density. Limits on substantive due process come from history and basic values, not an arbitrary line at the nuclear family. Brennan decried the ``cultural myopia'' of the nuclear-family model, stressing extended households' importance to Black and immigrant families. Burger would have required a variance application; Stewart and White found no fundamental right. Powell's test later anchored \textit{Glucksberg} and \textit{Dobbs}.

\textbf{Examples.} \emph{Same-Side:} prosecuting a woman for housing her aunt and sister. \emph{Opposite-Side:} a three-unrelated-roommates cap. \emph{Fence-Sitter:} a chosen family with an unadopted godchild.

\textbf{Critique.} Originalists say the test lets judges pick the level of generality; progressives praise its respect for kinship care but fault a blood-based definition of family that excludes chosen households.

\textbf{Key Quotation.} ``Our decisions establish that the Constitution protects the sanctity of the family precisely because the institution of the family is deeply rooted in this Nation's history and tradition.'' --- Justice Powell, plurality opinion, 431 U.S. at 503.

\end{document}
